% Zeros (v3 §2)
\begingroup\emergencystretch=2em % typography only: lets TeX loosen lines instead of overfilling them
\section{Zeros: everything reduces to a bound on the lag}\label{sec:zeros}

\begin{proposition}[The lag bound]\label{prop:La}
% v3: (L-a)
On the event $\mathcal G=\{\tau=\infty\}$ of Definition~\ref{def:good},
\begin{equation}\label{eq:La}
  |D(x)|\le x^{\eta}\qquad\text{for all }x\ge\Xm,\qquad \eta:=\min(\theta,\tfrac12-\theta)/2,
\end{equation}
and $D=0$ below $\Xm$.
\end{proposition}
\begin{proof}
% Logged correction 36: lag constant 4 on the lower side; start-up range re-derived with the exact accounting.
By \eqref{eq:lag} and Theorem~\ref{thm:queue}, $|D|\le B(x)+\frac{H(x)}2(1+\kappa_+/\log x)+4\le B(x)+H(x)+3$ for $x\ge X_0+2H$.
On $[\Xm,X_0+2H]$ we have $|D|\le B(X_0)+3H(X_0)+3$. Indeed, for $x\in[X_0,X_0+2H]$ the accounting in the proof of
Corollary~\ref{cor:lag} applies verbatim (with $j:=j_0-1$ if $x<s_{j_0}$, where $V_{j_0-1}=j_0$ and
$T(x)\in[j_0,j_0+\vartheta_{\mathrm s})$), and gives $|D|\le B(X_0)+\frac{H(X_0)}2(1+\kappa_+/\log X_0)+5\le B(X_0)+H(X_0)+5$. For
$x\in[\Xm,X_0)$, $D(x)$ is the number of jittered primes that land in $[\Xm,x]$, which is at most $H(X_0)/2+1$.

At every scale, $\log(B+3H+3)\le8\ell^2+10^{-6}L+2$. Indeed $\log H-\log H_{\mathrm{std}}\le10^{-6}L$ by
Theorem~\ref{thm:caplw}(c), and $B\le H$ by \eqref{eq:queueB}, since $\log M_1\le\frac12\log H$ and $4\ell^2\ge2\ell+3.6$. Since
$(8\ell^2+2)/L$ decreases for $L\ge\Lambda\ge10^8$, this is $\le2.82\cdot10^{-5}L$ for all $L\ge\Lambda-1$. For
$\theta\le\theta_1=\tfrac12-10^{-4}$ we have $\eta\ge5\cdot10^{-5}$, so the right side is $\le x^{\eta}$ for all $x\ge\Xm$.
\end{proof}

This is where $\theta\le\theta_1$ is used. At $X_0=e^{10^8}$ one has $\log H(X_0)=2715$, so \eqref{eq:La} fails near $X_0$ for
$\eta<2.8\cdot10^{-5}$; going beyond $\theta_1$ would need a $\theta$-dependent $X_0$, and one value of $\theta$ suffices for
Corollary~\ref{cor:notQ}. A polylogarithmic form of \eqref{eq:La} is false: the landing lag alone is $H/2=x^{o(1)}$.

\subsection{Decomposition of the zeta function}
For $\sigma>1$,
\begin{equation}\label{eq:zetadecomp}
  \zeta_{\PP}(s)=\int x^{-s}\,dN=Z_B(s)+\widehat E_{\Abs}(s)+\widehat R(s)+\widehat D(s),
\end{equation}
where:
\begin{itemize}
\item $Z_B(s):=\sum_{n<X_0}n^{-s}+\sum_{d\in\mathcal D,\,d\ge X_0}d^{-s}+\int_{X_0}^\infty x^{-s}(1-\rhodet)\,dx$;
\item $\widehat E_{\Abs}(s):=\int x^{-s}\,d\bigl(\#\Abs-\int\bar\rho_{\Abs}\bigr)$. On $\mathcal G$ the cap is inactive, and
$\widehat E_{\Abs}(s)=\sum_{a\in\Abs}a^{-s}(1-w_a(s))$, where
\[
  w_a(s):=a^s\int_{u_a}^au^{-s}g(u)\,du,\qquad g:=1/h,\qquad u_a+h(u_a)=a,
\]
is the weight of $a$ in the average $\rho_{\Abs}(u)=\#(\Abs\cap[u,u+h(u)])/h(u)$ over the forward window $[u,u+h(u)]$.
% Logged correction 33 (CORRECTIONS-vs-v3.md): v3's |w_a(s)| <= min(1, 2a^{1-theta/2}/|t|) is false (w_a(1) > 1); replaced by (i)-(iv).
Put $b:=\theta/2$, so that $h(u)=u^b$ and $g$ is decreasing, and $M:=\int_{u_a}^ag$. For $a\ge X_0/2$ (so for every $a\in\Abs$)
and $0\le\sigma\le2$:
\begin{enumerate}[label=(\roman*)]
\item $(u_a/a)^b\le M\le1$ and $a/u_a\le(1-a^{b-1})^{-1}$;
\item $|w_a(s)|\le(a/u_a)^{\sigma}M\le1+1.02\,\sigma a^{b-1}$;
\item $|w_a(s)|\le2.01\,a^{1-b}/|t|$ for $t\ne0$;
\item $|1-w_a(s)|\le1.01(|s|+b)a^{b-1}$; hence $|1-w_a(s)|\le2|t|a^{b-1}$ for $|t|\ge2$.
\end{enumerate}
Indeed, put $\epsilon:=a^{b-1}\le(X_0/2)^{-3/4}$.
(i) Since $g$ is decreasing, $h(u_a)/h(a)=(a-u_a)g(a)\le M\le(a-u_a)g(u_a)=1$, and $h(u_a)/h(a)=(u_a/a)^b$; also
$a-u_a=h(u_a)\le a^b$.
(ii) $|u^{-s}|\le u_a^{-\sigma}$ on $[u_a,a]$, and $(1-\epsilon)^{-\sigma}\le1+1.02\,\sigma\epsilon$.
(iii) Integrating, $w_a(s)=(a^{1-b}-a^su_a^{1-s-b})/(1-s-b)$, so $|w_a(s)|\le a^{1-b}\bigl(1+(a/u_a)^{\sigma+b-1}\bigr)/|t|$, and
$(a/u_a)^{\sigma+b-1}\le\max\bigl(1,(1-\epsilon)^{-2}\bigr)\le1.01$.
(iv) $1-w_a(s)=(1-M)+\int_{u_a}^a\bigl(1-(a/u)^s\bigr)g(u)\,du$. Here $1-M\le1-(1-\epsilon)^b\le1.01\,b\epsilon$, and, by
$|e^z-1|\le|z|e^{\max(0,\operatorname{Re}z)}$, $|1-(a/u)^s|\le|s|\log(a/u_a)(a/u_a)^{\sigma}\le|s|\epsilon(1-\epsilon)^{-3}\le1.01|s|\epsilon$;
with $M\le1$ this gives the first bound. For $|t|\ge2\ge\sigma$ we have $|s|\le\sqrt2|t|$, and $1.01(\sqrt2|t|+b)\le2|t|$.
\item $\widehat R:=\sum_k\widehat R_k$ with $\widehat R_k(s):=\int x^{-s}R_k'(x)\,dx$;
\item $\widehat D(s):=\int_{\Xm}^\infty x^{-s}\,dD=s\int_{\Xm}^\infty D(x)x^{-s-1}\,dx$, since $D=0$ below $\Xm$. Under \eqref{eq:La}
this converges absolutely for $\sigma>\eta$.
\end{itemize}

\begin{lemma}[The base factor]\label{lem:ZB}
% v3: Lemma 2.0
% Logged correction 67: explicit ranges (|t| >= Gamma_1/2; delta log X_0 <= 0.01) in place of "delta log X_0 -> 0".
Let $Z_\infty(s):=\prod_{p<X_0}(1-p^{-s})^{-1}=\sum_{n\in\mathcal D}n^{-s}$. For $|t|\ge\Gamma_1/2$ and $\tfrac12\le\sigma\le2$,
\[
  |Z_B(s)-Z_\infty(s)|\le\exp(-X_0^2/4).
\]
On $\sigma\ge1-\delta$ with $\delta\log X_0\le0.01$: $|Z_\infty|\le1.9\log X_0$ and $|Z_\infty'|\le4\log^2X_0$; if also
$|t|\ge\Gamma_1/2$ and $\sigma\le2$, then $|Z_B|\le2\log X_0=:Z_{\max}$.
\end{lemma}
\begin{proof}
% Logged correction 91 (CORRECTIONS-vs-v3.md): v3 gives the truncation as (X_0^{-s} - n^{-s})/(s kappa), which is not the omitted ramp
% (it vanishes at n = X_0); the omitted part is on [n/(1+kappa), X_0), with transform ((n/(1+kappa))^{-s} - X_0^{-s})/(s kappa).
Each rigid $n\ge X_0(1+\kappa)$ contributes a unit jump at $n$ and the ramp $-du/(\kappa u)$ on $[n/(1+\kappa),n]$. The at most
$\kappa X_0+1$ rigid $n\in[X_0,X_0(1+\kappa))$ have ramps truncated at $X_0$: the part of the ramp on $[n/(1+\kappa),X_0)$ is
omitted, which adds $\bigl((n/(1+\kappa))^{-s}-X_0^{-s}\bigr)/(s\kappa)$ to $Z_B$ compared with \eqref{eq:ZBformula}, of modulus
$\le2(1+\kappa)^{\sigma}X_0^{-\sigma}/(\kappa|s|)\le2.01X_0^{-\sigma}/(\kappa|s|)$ for $\sigma\le2$. So these boundary terms total
$\le2.01(X_0+1/\kappa)X_0^{-\sigma}/|s|\le3X_0^{1-\sigma}/|t|$. Up to them,
\begin{equation}\label{eq:ZBformula}
  Z_B=Z_\infty-W(s)S_{\ge X_0}(s)+\frac{X_0^{1-s}}{s-1},\qquad W(s):=\frac{(1+\kappa)^s-1}{s\kappa},\qquad
  S_{\ge X_0}(s):=\sum_{d\in\mathcal D,\,d\ge X_0}d^{-s}.
\end{equation}
For $\sigma\le2$, $|W|\le6/(\kappa|t|)$ and $|S_{\ge X_0}(s)|\le Z_\infty(\sigma)$. Also $\log Z_\infty(\frac12)\le3.42\sum_{2\le n<X_0}n^{-1/2}
\le6.9X_0^{1/2}$, as $-\log(1-y)\le y/(1-y)$, while $\kappa\Gamma_1\ge x_1^{1-\theta}/(2X_0^2\log^3x_1)\ge e^{0.49X_0^2}$, since
$\kappa,\kappa_0\ge X_0^{-1}$ and $\log x_1\ge X_0^2$; and the boundary terms and $X_0^{1-s}/(s-1)$ are $\le4X_0^{1/2}/|t|$ for
$\sigma\ge\frac12$. So for $|t|\ge\Gamma_1/2$ the difference is at most $(12e^{6.9X_0^{1/2}}+8X_0^{1/2})e^{-0.49X_0^2}\le e^{-X_0^2/4}$.
On $\sigma\ge1-\delta$ with $\delta\Lambda\le0.01$, we have $p^{\sigma}-1\ge0.98(p-1)$, so
$\sum_{p<X_0}\log p/(p^{\sigma}-1)\le1.03\Lambda$ by Mertens' estimates \cite{RosserSchoenfeld1962}. As
$(\log Z_\infty)'(\sigma)=-\sum_{p<X_0}\log p/(p^{\sigma}-1)$, this gives
$Z_\infty(\sigma)\le Z_\infty(1)e^{1.03\delta\Lambda}\le1.01e^{\gamma}\Lambda\cdot1.011\le1.9\Lambda$, and
$|Z_\infty'(s)|\le Z_\infty(\sigma)\sum_{p<X_0}\log p/(p^{\sigma}-1)\le4\Lambda^2$. Where also $|t|\ge\Gamma_1/2$,
$|Z_B|\le1.9\Lambda+e^{-X_0^2/4}\le2\Lambda$.
\end{proof}

\begin{lemma}[The base factor at all heights]\label{lem:ZBlow}
% v3: Lemma 2.0''
Let $0<a\le8$ and $\operatorname{Ein}(a):=\int_0^a(e^t-1)\,dt/t$, so that $\operatorname{Ein}(1)=1.32$ and $\operatorname{Ein}(8)=437.7$.
Put $\widehat E(s):=Z_B(s)-s/(s-1)$. Then for $\sigma\ge1-a/\Lambda$ and every $t$,
\[
\begin{aligned}
  |\widehat E(s)|&\le C_0e^a\Lambda+4.04e^{\operatorname{Ein}(a)}\min\bigl(1,\kappa(1+|s|)\bigr)e^{\gamma}\Lambda+3\kappa e^a+1\\
  &\le C_a(1+\kappa|s|)e^{\gamma}\log X_0,\qquad C_a:=C_0e^a+5e^{\operatorname{Ein}(a)},
\end{aligned}
\]
with $C_0:=2.1$. Also $1-W(1)=0$, so $Z_B$ has the pole and residue of $s/(s-1)$.
\end{lemma}
\begin{proof}
From \eqref{eq:ZBformula} and $s/(s-1)=1+\int_1^{X_0}x^{-s}\,dx+X_0^{1-s}/(s-1)$,
\[
  \widehat E(s)=\Bigl[\sum_{n<X_0}n^{-s}-1-\int_1^{X_0}x^{-s}\,dx\Bigr]+(1-W(s))S_{\ge X_0}(s)+(\text{boundary terms}).
\]
\emph{The bracket.} For $\sigma>2$ it is at most $2$; let $\sigma\le2$. For every integer $M\ge2$, partial summation gives the
exact identity
\begin{equation}\label{eq:zeta-partial}
  \zeta(s)-\frac1{s-1}=\sum_{n\le M}n^{-s}-\int_1^Mx^{-s}\,dx-s\int_M^\infty\{x\}x^{-s-1}\,dx .
\end{equation}
Let $|t|\le X_0$. With $M:=\lceil X_0\rceil-1$, \eqref{eq:zeta-partial} shows that the bracket equals $\zeta(s)-1/(s-1)-1$ up to
$s\int_M^\infty\{x\}x^{-s-1}\,dx-\int_M^{X_0}x^{-s}\,dx$, which has modulus $\le(|s|/\sigma+1)M^{-\sigma}\le1.1e^a$. (This makes
explicit the error term $O(X_0^{-\sigma})$ of \cite[Thm.~4.11]{Titchmarsh1986}.) With $M:=\lfloor|t|\rfloor+2$ instead, $|s|\le M+1$
and $\max(1,M^{1-\sigma})\le(X_0+2)^{a/\Lambda}\le1.001e^a$, so \eqref{eq:zeta-partial} gives
$|\zeta(s)-1/(s-1)|\le1.001e^a(1+2\log M+1.52)\le2.01e^a\Lambda$. For $|t|>X_0$ the bracket is
$\le2+2\int_1^{X_0}x^{a/\Lambda-1}\,dx\le2+2e^a\Lambda$. So the bracket is $\le C_0e^a\Lambda$ with $C_0:=2.1$.

\emph{The rigid part.} $|1-W(s)|\le\kappa|s-1|/2+O(\kappa^2(1+|s|)^2)\le\kappa(1+|s|)$ when $\kappa|s|\le1$, and $|W|\le3$ otherwise
(for $\sigma\le2$; for $\sigma>2$ the term is $\le2/X_0$). Next, $|S_{\ge X_0}(s)|\le Z_\infty(\sigma)$, and
\[
  \log Z_\infty(\sigma)-\log Z_\infty(1)\le\sum_{p<X_0}p^{-1}(p^{a/\Lambda}-1)+O(a/\Lambda)\le\operatorname{Ein}(a)+O(ae^a/\Lambda).
\]
So $Z_\infty(\sigma)\le1.01e^{\operatorname{Ein}(a)}e^{\gamma}\Lambda$ by Mertens' product.

% Logged correction 91: v3's '<= 3 n^{-sigma}' fails for large sigma ((1+kappa)^sigma); the omitted ramp lies on [n/(1+kappa),X_0), above 1.
\emph{Boundary.} The at most $\kappa X_0+1$ rigid $n\in[X_0,X_0(1+\kappa))$ each contribute the transform of the omitted ramp on
$[n/(1+\kappa),X_0)$, of modulus at most $\frac1\kappa\int_{n/(1+\kappa)}^{X_0}u^{-\sigma-1}\,du\le\frac{\log(1+\kappa)}\kappa
\bigl(\frac{X_0}{1+\kappa}\bigr)^{-\sigma}\le\bigl(\frac{X_0}{1+\kappa}\bigr)^{a/\Lambda-1}\le1.01X_0^{a/\Lambda-1}$: the support lies
above $1$, so the bound decreases in $\sigma$, and $\sigma\ge1-a/\Lambda$. In total this is
$\le1.01(\kappa X_0+1)e^aX_0^{-1}\le3\kappa e^a$.
\end{proof}

The local prime number theorem (Theorem~\ref{thm:lp}) compares $N$ with the target without its absorbed part,
$T_0:=T-(\#\Abs-\int\bar\rho_{\Abs})$, whose transform is $\widehat T_0=Z_B+\widehat R=s/(s-1)+\widehat R+\widehat E$. The absorbed part
$\widehat E_{\Abs}$ is \emph{not} bounded on all of $\sigma\ge1-a/\Lambda$ at large $|t|$ (Remark~\ref{rem:EAhuge}), so it is not put
% Logged correction 34: the constant 2 became 3.1 (mass deficit of the smoothed absorbed elements).
into $\widehat T_0$. On $\Good(<y)$, $|N-T_0|\le|N-T|+3.1u^{\theta/2}(\log u)^{-4}+2$ on $[1,y]$ by the absorbed part of
\eqref{eq:T1}. The constant of Lemma~\ref{lem:ZBlow} enters only through $\log C_a$; with $C_0=2.1$, $\log C_8<439.4$.

\begin{lemma}[The absorbed part]\label{lem:ZA}
% v3: Lemma 2.0'
Put $u_*:=\Lambda e^{s_H}/22.6$, with $s_H=5.9745$ as in Theorem~\ref{thm:caplw}(b); so $u_*\ge1.7\cdot10^9$.
\begin{enumerate}[label=(\alph*)]
% Logged correction 35 (CORRECTIONS-vs-v3.md): v3 stated 2exp(-u0 log u0/2); the proof gives the bound below
% (with the constant u_* in place of u_0, since the saving is not monotone across the ramp; R4 N12).
\item On $\mathcal G$,
$|\widehat E_{\Abs}(s)|\le2\exp(-0.45\,u_*\log u_*)+5.4(\log x_k)^9\exp\bigl(-\log x_k/(22\log\log x_k)\bigr)\le e^{-10^{10}}$,
where $k$ is the index of the designed zero concerned. This holds on every disc
$|s-(1+i\gamma_k)|\le\delta_k$, on every Rouch\'e circle of Lemma~\ref{lem:designedzero}, and on $\sigma=1$ at the designed heights.
\item On $\mathcal G$, the term $\int_{X_0}^{x_k}x^{-s}\rho_{\Abs}\,dx$ has mean square $\le x_k^{\theta-1/2-o(1)}$ over
$t\in[\Gamma_k,2\Gamma_k]$, $\sigma\in[\tfrac34,\tfrac54]$.
\end{enumerate}
\end{lemma}
\begin{proof}
The elements of $\Abs$ are $2y$-smooth, and Rankin's method for the realised system (Lemma~\ref{lem:lemmaP}, proof steps (i) and (v),
with $Y=x$) gives $\#\Abs\cap[1,x]\le x\exp(-0.9u\log u)$, with $u=\log x/\log(2y)$.
% Logged correction 65 (CORRECTIONS-vs-v3.md): step (v) of Lemma P now carries the hypothesis (*); "u >= 100" was too weak for it.
If $\Abs\cap[1,x]\ne\emptyset$, the hypothesis $(*)$ of step (v) holds exactly as in step (vi) there: $u\ge u_0\ge2\Lambda$, with $u_0$
as in \eqref{eq:twelvetwo} at height $x$, $\log\log2y\le1.79u$, and $\log u/\log2y\le10^{-4}$ (the $\eta$ of step (v)); otherwise there is nothing to prove. Moreover
$\Abs=\emptyset$ wherever $y<\Xm$, which holds for $\log L\le2\log\Lambda+s_H$ (Theorem~\ref{thm:caplw}(b)). Where $\Abs\ne\emptyset$ and $\ell\le\ell_{\mathrm{std}}$,
$\log2y\le22.6\Lambda$ (Section~\ref{sec:params}). In the standard regime $\ell\ge\ell_{\mathrm{std}}$ instead $H=H_{\mathrm{std}}$,
$c=\tfrac14H_{\mathrm{std}}^{1/2}$ and $\log2y\le12\ell^2+9\ell+8\le12.1\ell^2$.
% R4 N12: an increasing minorant of the saving replaces v3's "rate >= 0.18/Lambda" (u = v/log 2y is not monotone across the ramp).
So, in the variable $v=\log x$, wherever $\Abs\ne\emptyset$ we have $\log2y\le\lambda(v):=\max\bigl(22.6\Lambda,\,12(\log v)^2+9\log v+8\bigr)$,
and therefore
\[
  \#\Abs\cap[1,x]\le x\,e^{-\varsigma_*(\log x)},\qquad \varsigma_*(v):=0.9\,\tilde u(v)\log\tilde u(v),\quad \tilde u(v):=v/\lambda(v).
\]
The minorant $\varsigma_*$ is continuous and increasing, since $v\lambda'(v)\le24\log v+9<\lambda(v)$. For $\log v\ge\ell_{\mathrm{std}}$ we have
$\lambda(v)\le12.1(\log v)^2$, hence $\varsigma_*(v)\ge v/(14\log v)$. Since $\Abs=\emptyset$ for $\log v\le2\log\Lambda+s_H$, the
least $v_0$ with $e^{v_0}\in\Abs$ satisfies $\tilde u(v_0)\ge\Lambda^2e^{s_H}/(22.6\Lambda)=u_*$, so
$\varsigma_*(v_0)/2\ge0.45\,u_*\log u_*\ge1.66\cdot10^{10}$ for every $\Lambda\ge10^8$.

(a) Fix $|t|\in[\Gamma_k/2,3\Gamma_k]$ and $1-\delta_k\le\sigma\le2$ (this covers all the sets in (a)), and put
$Y:=|t|^{1/(1-\theta/2)}\le x_k$.
% Logged correction 33: |1-w_a| <= 2.0001 (from (ii)) replaces |1-w_a| <= 2; the a > Y bound is (iv), now proved.
For the terms $a\le Y$ we use $|1-w_a(s)|\le1+|w_a(s)|\le2+2.1a^{\theta/2-1}\le2.0001$, by (ii) after \eqref{eq:zetadecomp}. By
partial summation these terms are then at most $2.0001$ times
$\exp(\delta_k\log Y-\varsigma_*(\log Y))+\int_{v_0}^{\log Y}\exp(\delta_kv-\varsigma_*(v))\,dv$. This is
$\le2\exp(-\varsigma_*(v_0)/2)\le2\exp(-0.45\,u_*\log u_*)$, uniformly in $k$.
% R4 N5: the bound is made uniform in k (v3: "the length of the range is absorbed").
Indeed, $\varsigma_*$ is increasing, so $e^{-\varsigma_*(v)}\le e^{-\varsigma_*(v_0)/2}e^{-\varsigma_*(v)/2}$, and it suffices that
$e^{\delta_k\log Y-\varsigma_*(\log Y)/2}+\int_{v_0}^{\log Y}e^{\delta_kv-\varsigma_*(v)/2}\,dv\le10^{-6}$. Put $V:=e^{\ell_{\mathrm{std}}}$. On
$[v_0,V]$ we have $\delta_kv\le18\Lambda e^{-2\Lambda}V\le10^{-6}$, since $\delta_k\le9\log\log x_k/\log x_k\le18\Lambda e^{-2\Lambda}$,
and $\varsigma_*(v)/2\ge\varsigma_*(v_0)/2\ge\ell_{\mathrm{std}}+20$, so that part is $\le1.01\,V e^{-\ell_{\mathrm{std}}-20}\le10^{-8}$. On
$[V,\log Y]$ we have $\log v\le\log\log x_k$, hence $\delta_kv\le\varsigma_*(v)/4$ (as $504(\log\log x_k)^2\le\log x_k$), and
$\varsigma_*(v)\ge v/(14\log v)$; so that part, and the boundary term, are $\le e^{-V/(60\log V)}\le10^{-8}$. The terms $a>Y$ use $|1-w_a|\le2|t|a^{\theta/2-1}$, from (iv) after
\eqref{eq:zetadecomp}. With $b=\theta/2$, $|t|=Y^{1-b}$ and partial summation, they are at most
\[
  2|t|\sum_{a\in\Abs,\,a>Y}a^{b-2+\delta_k}\le2Y^{1-b}\,\frac{2-b}{1-b-\delta_k}\,Y^{b-1+\delta_k}e^{-\varsigma_*(\log Y)}\le5.4\,Y^{\delta_k}e^{-\varsigma_*(\log Y)} .
\]
Here $Y\le x_k$, so $Y^{\delta_k}\le x_k^{\delta_k}\le(\log x_k)^9$. Also $\log Y\ge0.66\log x_k$, since
% Logged correction 67: explicit lower bound for log Gamma_k (v3: "- O(log log x_k)").
$\log\Gamma_k\ge(1-\theta)\log x_k-4\log\log x_k$ (as $\log(1/\kappa_0)\le\Lambda\le\frac12\log\log x_k$) and
$(1-\theta)/(1-\theta/2)\ge\tfrac23+8\cdot10^{-5}$. These $a$ lie in the standard regime
($\log\log Y\ge2\Lambda-1>\ell_{\mathrm{std}}$), where $\varsigma_*(\log Y)\ge\log Y/(14\log\log Y)\ge\log x_k/(22\log\log x_k)$.
So these terms are $\le5.4(\log x_k)^9\exp(-\log x_k/(22\log\log x_k))$, which is $\le e^{-10^{11}}$ since $\log x_k\ge X_0^2$.

(b) Mellin--Plancherel gives, at $\sigma=\tfrac34$,
$(4\pi/\Gamma_k)\sum_{a\in\Abs,\,a\le2x_k}a^{1-2\sigma}\le x_k^{1/2-1/(14\log\log x_k)}/\Gamma_k+2\exp(e^{\ell_{\mathrm{std}}}/2)/\Gamma_k$,
and $\Gamma_k\asymp x_k^{1-\theta}/\log^3x_k$.
% Logged correction 67: part (b) is no longer used (correction 64); its o(1) is therefore not load-bearing.
Part (b) is not used for Theorem~\ref{thm:A}: the proof of Lemma~\ref{lem:goodheights}(ii) bounds this term directly, with explicit
constants.
\end{proof}

\begin{remark}\label{rem:EAhuge}
$\widehat E_{\Abs}$ can be huge off these sets, for example at $\sigma=1-1/\log X_0$. No step of this section evaluates it there.
Lemma~\ref{lem:ZA} needs only $y\le n^{o(1)}$ on the range of $\Abs$, $\log y\le23\log X_0$ for $\ell\le\ell_{\mathrm{std}}$ and
$\log2y\le12.1\ell^2$ beyond, so it holds for every absorption threshold of Section~\ref{sec:params}, uniformly in $\Lambda\ge10^8$.
\end{remark}

\subsection{The oscillations}
Put $\Phi(w):=\int_0^1e^{-wv}\phi'(v)\,dv$. Then
\begin{equation}\label{eq:Rhat}
  \widehat R_k(s)=\frac{\alpha_ks}2\Bigl[x_k^{-w}\frac{\Phi(w)}w+x_k^{-w'}\frac{\Phi(w')}{w'}\Bigr],\qquad w=s-\theta-i\gamma_k,\quad
  w'=s-\theta+i\gamma_k.
\end{equation}
\begin{itemize}
% Logged correction 67 (CORRECTIONS-vs-v3.md): the decay constants C_m, C_2 and the o(1) bounds near gamma_k are replaced by
% explicit bounds that use only |Phi| <= 1.
\item $\Phi$ is entire, with $\Phi(0)=1$, and $|\Phi(w)|\le1$ and $|\Phi'(w)|\le1$ for $\operatorname{Re}w\ge0$, as $\phi'\ge0$ and
$\int_0^1\phi'=1$.
\item $\operatorname{Re}\Phi(w)\ge e^{-1}\cos1=:\operatorname{Re}\Phi_{\min}$ for $0\le\operatorname{Re}w\le1$ and $|\operatorname{Im}w|\le1$.
\item So $\zeta_{\PP}$ continues to $\sigma>\max(\theta,\eta)$, except for the simple pole at $1$.
\item Write $M_k$ for the resonant term (the one with $w$) and $E_k$ for the conjugate one. On $0.99\le\sigma\le2$, $|t-\gamma_k|\le1$,
\[
  |E_k|\le0.51\,x_k^{\theta-\sigma},\qquad\sum_{j<k}|\widehat R_j|\le2.02\,x_1^{\theta-\sigma},\qquad
  \sum_{j>k}|\widehat R_j|\le4.2\,x_ke^{-0.24x_k}.
\]
Indeed $\operatorname{Re}w,\operatorname{Re}w'\ge0.49$ there and $|\Phi|\le1$. For $E_k$, $|s|\le0.5001|w'|$. For $j<k$, $|s|\le1.0001|w|$
and $|s|\le|w'|$, since $\gamma_j\le x_j\le2\log x_k$ is negligible against $\gamma_k\ge\Gamma_k$; and the $x_j$ grow
super-exponentially. For $j>k$, $|w|,|w'|\ge0.49$, $|s|\le\gamma_k+3\le x_k$ and $x_j^{\theta-\sigma}\le x_j^{-0.49}$ with
$x_j\ge\exp(x_k/2)$ (the spacing $x_j\ge\exp(x_{j-1})$ of Section~\ref{sec:params} holds for the construction scales). Together
the three bounds are at most $3e^{-0.49X_0^2}$.
% Logged correction 76 (R17, C76-3): with x_j = x_j^+ here, only x_j >= exp(x_k/2) holds for j > k; the exponent 0.49x_k
% becomes 0.24x_k, and the total 3e^{-0.49 X_0^2} is unaffected.
\item On $\sigma\ge1-8/\log X_0$, at every height,
\begin{equation}\label{eq:Rhatbound}
  |\widehat R(s)|\le|s|\,e^{-0.39X_0^2}.
\end{equation}
Indeed there $\operatorname{Re}w,\operatorname{Re}w'\ge\tfrac12-8/\log X_0>0.4$, $|\Phi|\le\int\phi'=1$ and $|1/w|\le2.5$, so
$|\widehat R_k|\le5|s|x_k^{-0.4}$. Theorem~\ref{thm:lp} uses this.
\end{itemize}

\subsection{The designed zeros}
\begin{lemma}[A designed zero]\label{lem:designedzero}
% v3: Lemma 2.1
% R4 N4: the constant in delta_k (v3: C_1, unspecified) is set to 0; step 6 of the proof checks that this is admissible.
Suppose $|Z_B|\in[\tfrac12,Z_{\max}]$ on $|s-(1+i\gamma_k)|\le\delta_k$, where $\delta_k:=9\log\log x_k/\log x_k$. Then
$Z_B+M_k$ has a zero $s^*$ with
\[
  \operatorname{Re}s^*=1-\frac{3\log\log x_k+O_{X_0}(1)}{\log x_k},
\]
and $|Z_B+M_k|\ge\tfrac14$ on the circle $|s-s^*|=1/\log x_k$.
\end{lemma}
\begin{proof}
\begin{enumerate}[label=\arabic*.]
\item The equation is $x_k^w=G(w):=-(\alpha_ks/2)\Phi(w)/(wZ_B(s))$.
\item Since $\alpha_k\gamma_k=\kappa_0x_k^{1-\theta}/\log^3x_k$, we have $\log|G|=(1-\theta)\log x_k-3\log\log x_k+O_{X_0}(1)$ on the disc $|s-(1+i\gamma_k)|\le\delta_k$.
% Logged correction 67: the o(1) terms of steps 3-5 are written out with their constants.
\item On the disc, $|Z_B'|\le4.01\Lambda^2$: Lemma~\ref{lem:ZB} gives $|Z_\infty'|\le4\Lambda^2$, and Cauchy's estimate on discs of radius
$\frac14$ gives $|(Z_B-Z_\infty)'|\le4e^{-X_0^2/4}$. With $|\Phi'|\le1\le e|\Phi|$ and $|w|\ge0.49$ (step~6), and $|s|\ge\gamma_k-1\ge10^6$,
\[
  \Bigl|\frac{G'}{G}\Bigr|=\Bigl|\frac1s+\frac{\Phi'}{\Phi}-\frac1w-\frac{Z_B'}{Z_B}\Bigr|\le4.8+8.02\Lambda^2\le8.1\Lambda^2.
\]
Fix the branch of $\log G$ with $|\operatorname{Im}\log G|\le\pi$ at the centre of the disc; then $|\operatorname{Im}\log G|\le\pi+0.01$
on the disc. The map $F(w):=(\log G(w))/\log x_k$ has $|F'|\le8.1\Lambda^2/\log x_k\le8.1\Lambda^2/X_0^2<1$, and by the bounds of
step~6, which hold at every point of the disc, it maps the disc (in the variable $w$) into itself. So it has a unique fixed point
$w^*$, and $\operatorname{Re}s^*=\theta+\operatorname{Re}w^*$.
\item On the circle, with $\xi:=(w-w^*)\log x_k$: $M_k(s)/M_k(s^*)=e^{-\xi}(1+\epsilon_1)$ and $Z_B(s)/Z_B(s^*)=1+\epsilon_2$, where
$|\epsilon_1|\le5/\log x_k$, since $\log(s\Phi(w)/w)$ has derivative at most $10^{-6}+e+2.05\le4.8$, and
$|\epsilon_2|\le8.1\Lambda^2/\log x_k$, by step~3 and $|Z_B(s^*)|\ge\frac12$.
\item Since $Z_B(s^*)=-M_k(s^*)$, we get $Z_B+M_k=Z_B(s^*)\bigl(1-e^{-\xi}+\epsilon_2-e^{-\xi}\epsilon_1\bigr)$. On $|\xi|=1$ we have
$|1-e^{-\xi}|\ge1-e^{-1}$ (the minimum is at $\xi=1$) and $|e^{-\xi}|\le e$, so
$|Z_B+M_k|\ge\frac12(1-e^{-1}-8.1\Lambda^2/X_0^2-5e/X_0^2)\ge0.31\ge\tfrac14$.
\item The circle lies in the disc. In step~2 the $O_{X_0}(1)$ term lies in $[\log\kappa_0-\log(2\Lambda)-1.5,\ \log\kappa_0+0.72]$,
since $e^{-1}\le|\Phi(w)|\le1$, $0.49\le|w|\le0.76$ and $\tfrac12\le|Z_B|\le2\Lambda$ there; it is negative, as $\kappa_0\le\tfrac14$.
With the branch of step~3, $|\operatorname{Im}F|\le(\pi+0.01)/\log x_k$. Hence
$|s^*-(1+i\gamma_k)|+1/\log x_k\le(3\log\log x_k+\log(1/\kappa_0)+\log(2\Lambda)+5.7)/\log x_k\le\delta_k$, because
$\kappa_0\ge X_0^{-1}$ and $\log\log x_k\ge2\Lambda$.\qedhere
\end{enumerate}
\end{proof}

The discs $|s-(1+i\gamma_k)|\le\delta_k$ are small for every $k\ge1$: $\delta_k\le\delta_1\le18\log X_0/X_0^2$, so $\delta_k\log X_0\le0.01$
and $|t|\ge\Gamma_k-1\ge\Gamma_1/2$ there. So $Z_B$ is evaluated only where Lemma~\ref{lem:ZB} applies.

\begin{lemma}[Late noise]\label{lem:latenoise}
% v3: Lemma 2.2
% Logged correction 67: the limit "-> 0" is replaced by the explicit bound 165 Lambda x_k^{-1/8}.
% Logged correction 76 (CORRECTIONS-vs-v3.md): the tail now starts at the construction scale a_k (Lemma goodheights), so it
% contains the strip [a_k, x_k]; the comparison with s^* (built from x_k = x_k^+) is redone with a_k.
On the circle of Lemma~\ref{lem:designedzero}, with $a_k$ as in Lemma~\ref{lem:goodheights} and
$\widehat D_{\ge a_k}(s):=\int_{[a_k,\infty)}x^{-s}\,dD$,
\[
  |\widehat D_{\ge a_k}(s)|\le6.001e\cdot|w^*|\cdot x_k^{\eta-\theta}\cdot\frac{Z_{\max}}{\alpha_k\operatorname{Re}\Phi_{\min}}
  \le165\,\Lambda\,x_k^{-1/8}\qquad(|w^*|\le1).
\]
\end{lemma}
\begin{proof}
By parts, $|\int_{[a_k,\infty)}x^{-s}\,dD|\le3|s|a_k^{\eta-\sigma}$, using \eqref{eq:La}. On the circle $\sigma\le1.01$, so
$a_k^{\eta-\sigma}\le x_k^{\eta-\sigma}(x_k/a_k)^{\sigma}\le(1+5H(a_k)/a_k)\,x_k^{\eta-\sigma}$, as $x_k=a_k+2H(a_k)$. Also
$x_k^{-\sigma}\le e\,x_k^{-\operatorname{Re}s^*}=e\,x_k^{-\theta}/|G(w^*)|$, and $|G(w^*)|\ge|s^*|\alpha_k\operatorname{Re}\Phi_{\min}/(2|w^*|Z_{\max})$.
The factors $|s|/|s^*|\le1.0001$ and $1+5H(a_k)/a_k\le1+10^{-10}$ cost at most $6\cdot1.0001\cdot(1+10^{-10})\le6.001$ in place of
$6$. For the last bound, $|w^*|\le1$,
$Z_{\max}=2\Lambda$, $\alpha_k\ge1$, $\operatorname{Re}\Phi_{\min}=e^{-1}\cos1\ge0.198$, and $\eta-\theta\le-\theta/2\le-\frac18$.
\end{proof}

\begin{lemma}[Good heights exist]\label{lem:goodheights}
% v3: Lemma 2.3
% Logged correction 76 (CORRECTIONS-vs-v3.md): v3 reads D up to x_k, which in this section means x_k^+ = a_k + 2H(a_k); the
% system on [1,x_k^+) involves slots past x_k^+, where R_k, and so gamma_k, enters T. The information cutoff is now the
% construction scale a_k, and gamma_k is defined for every omega by an explicit measurable rule (v3 only asserts existence).
Let $\Gamma_k:=\kappa_0x_k^{1-\theta}/(2\log^3x_k)$, and let $a_k$ be the construction scale of Section~\ref{sec:target}, so that
$x_k=a_k+2H(a_k)$ in this section. Put $\widehat D_{<a_k}(s):=\int_{[\Xm,a_k)}x^{-s}\,dD$, and let $\mathcal S_k$ be the set of
$\gamma\in[\Gamma_k,2\Gamma_k]$ with:
\begin{enumerate}[label=(\roman*)]
\item $|Z_B(1+i\gamma)|\ge0.6$;
\item $|\widehat D_{<a_k}(s)|\le1/16$ on $|s-(1+i\gamma)|\le\tfrac18$.
\end{enumerate}
Put $\gamma_k:=\min\mathcal S_k$ if $\mathcal S_k\ne\emptyset$, and $\gamma_k:=\Gamma_k$ otherwise. Then $\gamma_k\in[\Gamma_k,2\Gamma_k]$ is
defined for every $\omega$, and it is a measurable function of the system on $[1,a_k)$, which is decided at slots
$\le a_k+0.51H(a_k)<x_k$ and so does not depend on $\gamma_k$. So $R_k$, and hence $T$, is adapted (Lemma~\ref{lem:adapted}).
On $\Good(<a_k+H(a_k))$, which contains $\mathcal G$, the set $\mathcal S_k$ has measure at least $0.018\,\Gamma_k$; in particular
$\gamma_k\in\mathcal S_k$ there.
\end{lemma}
\begin{proof}
% Logged correction 64 (CORRECTIONS-vs-v3.md): v3 used a mean value theorem with an unspecified constant and concluded that the
% bad set is small "once X_0 is large"; every constant below is explicit, so the lemma holds for every admissible X_0.
% Logged correction 76 (R17, C76-5): the bound a_k + 0.51H(a_k) on the slots also uses eq:RH.
\emph{The rule.} Both conditions are closed in $\gamma$: $Z_B$ is continuous on $\sigma=1$, $t\ne0$ \eqref{eq:ZBformula}, and
$\widehat D_{<a_k}$ is entire, so $g(\gamma):=\max_{|s-(1+i\gamma)|\le1/8}|\widehat D_{<a_k}(s)|$ is continuous. Hence $\mathcal S_k$ is
compact, and $\min\mathcal S_k$ exists if $\mathcal S_k\ne\emptyset$. Put $h(\gamma):=\max(0.6-|Z_B(1+i\gamma)|,\,g(\gamma)-\frac1{16})$, a
continuous function of $\gamma$ and a measurable function of the system on $[1,a_k)$ (which determines the finite measure $dD$ on
$[\Xm,a_k)$). For $c\in[\Gamma_k,2\Gamma_k)$, $\{\gamma_k\le c\}=\{\min_{[\Gamma_k,c]}h\le0\}\cup\{\min_{[\Gamma_k,2\Gamma_k]}h>0\}$, and each
minimum is an infimum over a countable dense set; so $\gamma_k$ is measurable. The system on $[1,a_k)$ consists of the elements
below $a_k$, whose primes are decided at slots $\le a_k+0.51H(a_k)$ (as $s-H(s)/2$ increases, and $H(a_k+0.51H(a_k))\le1.01H(a_k)$
by \eqref{eq:RH}), and of $T$ on $[1,a_k)$; those slots
and their decisions are read off $T$ on $[1,a_k+H(a_k)]$, where $R_k\equiv0$ (\S\ref{sec:target}). The fallback $\gamma_k=\Gamma_k$
gives $\alpha_k=2$ in Theorem~\ref{thm:zeros}, so $\alpha_k\in[1,2]$ and \eqref{eq:ampl}, hence \eqref{eq:G1}, hold for every
$\omega$.

Put $\Gamma:=\Gamma_k$. The bounds on $\mathcal S_k$ use the mean value theorem of Montgomery and Vaughan. Let $\mu_r$ be distinct reals,
$\delta_r:=\min_{r'\ne r}|\mu_r-\mu_{r'}|$, and $b_r$ coefficients with $\sum_r|b_r|<\infty$ and $\sum_r|b_r|^2/\delta_r<\infty$. Then, for
all $T_0\in\R$ and $T>0$,
\begin{equation}\label{eq:MV}
  \Bigl|\int_{T_0}^{T_0+T}\Bigl|\sum_rb_re^{i\mu_rt}\Bigr|^2dt-T\sum_r|b_r|^2\Bigr|\le3\pi\sum_r\frac{|b_r|^2}{\delta_r}.
\end{equation}
For finite sums, expanding the square leaves the difference of two sums $\sum_{r\ne r'}c_r\bar c_{r'}/(i(\mu_r-\mu_{r'}))$ with
$|c_r|=|b_r|$, and each is at most $\frac{3\pi}2\sum_r|b_r|^2/\delta_r$ by the weighted Hilbert inequality of
\cite{MontgomeryVaughan1974}. Infinite sums follow by truncation, which does not decrease the $\delta_r$.

\emph{(i)} Apply \eqref{eq:MV} with $\mu_n=-\log n$, $n\in\mathcal D$, for which $\delta_n\ge\log(1+1/n)\ge1/(2n)$, to $Z_\infty(1+it)$
and to $Z_\infty(1+it)^2=\sum_{n\in\mathcal D}\tau(n)n^{-1-it}$. Over $t\in[\Gamma,2\Gamma]$,
\[
\begin{gathered}
  \Bigl|\E|Z_\infty(1+it)|^2-\prod_{p<X_0}(1-p^{-2})^{-1}\Bigr|\le\frac{6\pi Z_\infty(1)}\Gamma,\\
  \E|Z_\infty|^4\le\sum_{n\in\mathcal D}\frac{\tau(n)^2}{n^2}+\frac{6\pi}\Gamma\sum_{n\in\mathcal D}\frac{\tau(n)^2}n
  \le\frac{\zeta(2)^4}{\zeta(4)}+\frac{6\pi Z_\infty(1)^4}\Gamma\le6.8.
\end{gathered}
\]
Here $\sum_{n\in\mathcal D}\tau(n)^2/n=\prod_{p<X_0}(1+p^{-1})(1-p^{-1})^{-3}\le Z_\infty(1)^4$, $Z_\infty(1)\le1.01e^{\gamma}\Lambda$ by
Mertens' theorem, $\zeta(2)^4/\zeta(4)=5\pi^4/72<6.77$, and $\Gamma\ge e^{X_0^2/3}$ by \eqref{eq:Gammalow} below; so the error terms are
at most $200\Lambda^4e^{-X_0^2/3}$, and $\E|Z_\infty(1+it)|^2\in[1,\zeta(2)]$. The Paley--Zygmund inequality gives
$|Z_\infty(1+it)|\ge2^{-1/2}$ on a set of $t$ of density $\ge\frac14/6.8>0.036$, independently of $X_0$. Lemma~\ref{lem:ZB} moves this by
at most $\exp(-X_0^2/4)$.

% Logged correction 73 (CORRECTIONS-vs-v3.md): v3 works on Good(<x_k), where the queue bound, hence eq:La, is certified only
% up to x_k - 12 log x_k - 3; the cutoff is raised to a_k + H(a_k). This changes no event time: the choice of gamma_k still
% reads only the system on [1,a_k), and only its existence on G is used (Theorem thm:zeros).
% Logged correction 76: the proof of (ii) runs with a_k in place of x_k^+; Gamma and the disc are unchanged.
% Logged correction 76 (R17, C76-2): gamma_j <= x_j <= 2V, as x_{k-1} = x_{k-1}^+ <= 2a_{k-1} <= 2V.
\emph{(ii)} Write $f:=\widehat D_{<a_k}$ and $V:=\log a_k\ge X_0^2$, and work on $\Good(<a_k+H(a_k))$, where \eqref{eq:La} holds on
$[\Xm,a_k+2]$ (Remark~\ref{rem:horizon}). This event contains $\mathcal G$, on which Theorem~\ref{thm:zeros} uses (ii).
Since $\kappa_0\ge X_0^{-1}$, and $u^{1-\theta}/\log^3u$ increases for $\log u\ge6$ while $x_k\ge a_k$,
\begin{equation}\label{eq:Gammalow}
  \Gamma\ge\frac{a_k^{1-\theta}}{2X_0V^3},\qquad 1-\theta\ge\tfrac12+10^{-4}.
\end{equation}
Let $Q:=[\frac34,\frac54]\times[\Gamma-\frac14,2\Gamma+\frac14]$ in the coordinates $(\sigma,t)$, and
$\|F\|:=(\Gamma^{-1}\iint_Q|F|^2\,d\sigma\,dt)^{1/2}$. As the area of $Q$ is at most $\Gamma$, $\|F\|\le\sup_Q|F|$.
\begin{itemize}
\item \emph{Reduction.} The function $f$ is entire. If $|s-(1+i\gamma)|\le\frac18$, the disc $|z-s|\le\frac18$ lies in
$Q_\gamma:=[\frac34,\frac54]\times[\gamma-\frac14,\gamma+\frac14]$, so by subharmonicity $|f(s)|^2\le(64/\pi)\iint_{Q_\gamma}|f|^2$. If (ii)
fails at $\gamma$, then $\iint_{Q_\gamma}|f|^2>\pi/16384$. Each point of $Q$ lies in $Q_\gamma$ for a set of $\gamma$ of measure
$\le\frac12$, so the $\gamma\in[\Gamma,2\Gamma]$ at which (ii) fails have measure $\le(8192/\pi)\,\Gamma\|f\|^2$. We show
$\|f\|\le12X_0^{-3/16}$; then this measure is $\le3.8\cdot10^5X_0^{-3/8}\,\Gamma<0.018\,\Gamma$.
\item \emph{Decomposition.} With $\mathcal B:=\NN\setminus(\mathcal D\cup\Abs)$, counted with multiplicity,
\[
  f=\sum_{n\in\mathcal B,\ \Xm\le n<a_k}n^{-s}-\int_{X_0}^{a_k}x^{-s}(1-\rhodet-\bar\rho_{\Abs}+R')\,dx,
\]
by \eqref{eq:target}, and by Minkowski's inequality it suffices to bound the norms of the five pieces.
\item \emph{The $1$ and $\rhodet$ parts.} On $Q$, $|\int_{X_0}^{a_k}x^{-s}dx|\le2a_k^{1/4}/|t|$. The $\rhodet$ part is a truncated
$\sum_{d\in\mathcal D}d^{-s}W(s)$, with $W$ as in \eqref{eq:ZBformula}. Each rigid $d$ with $[d/(1+\kappa),d]\subset[X_0,a_k]$ contributes
$d^{-s}W(s)$, and $|W(s)|\le3/(\kappa|t|)$. The at most $\kappa X_0+1$ rigid $d\in[X_0,(1+\kappa)X_0)$, and the at most $\kappa a_k+1$
rigid $d\in(a_k,(1+\kappa)a_k]$, have truncated ramps, each of modulus $\le2.01X_0^{-\sigma}/(\kappa|t|)$, respectively
$\le2.01(a_k/(1+\kappa))^{-\sigma}/(\kappa|t|)$. So this part is at most $(3Z_\infty(\frac34)/\kappa+3X_0^{1/4}+3a_k^{1/4})/|t|$, where
$\log Z_\infty(\frac34)\le2.47\sum_{2\le n<X_0}n^{-3/4}\le10X_0^{1/4}$ and $1/\kappa\le X_0$. As $|t|\ge0.99\Gamma$, the two parts together
are at most $6a_k^{1/4}/\Gamma\le12X_0V^3a_k^{-1/4}$ on $Q$, by \eqref{eq:Gammalow}.
\item \emph{The $R'$ part.} On $[X_0,a_k]$ we have $R'=\sum_{j<k}R'_j$, and
$\int_{X_0}^{a_k}x^{-s}R'_j\,dx=\widehat R_j(s)-\int_{a_k}^\infty x^{-s}R'_j\,dx$. On $Q$, $\operatorname{Re}w,\operatorname{Re}w'\ge\frac14$ in
\eqref{eq:Rhat}, so $|\Phi|\le1$ there; and $|s|\le1.0001\min(|w|,|w'|)$, since $\gamma_j\le x_j\le2V$ is negligible against $\Gamma$. So
$|\widehat R_j(s)|\le2.01x_j^{-1/4}$. Beyond $a_k$ we have $\phi=1$ in $R_j$, so $|R'_j(x)|\le2(\gamma_j+1)x^{\theta-1}$, and the tail
is at most $8(\gamma_j+1)a_k^{-1/4}$. With $\sum_{j<k}(\gamma_j+1)\le12V$, the $R'$ part is at most
$2.02x_1^{-1/4}+100Va_k^{-1/4}\le2.02e^{-X_0^2/4}+100Va_k^{-1/4}$ on $Q$.
\item \emph{The $\bar\rho_{\Abs}$ part.} By Plancherel in $\log x$, for each $\sigma$,
$\int_\R|\int_{X_0}^{a_k}x^{-\sigma-it}\bar\rho_{\Abs}(x)\,dx|^2dt=2\pi\int_{X_0}^{a_k}x^{1-2\sigma}\bar\rho_{\Abs}(x)^2\,dx$. Here
$\bar\rho_{\Abs}\le0.1/\log x\le0.1/\Lambda$ by the cap in \eqref{eq:target}, and $x^{1-2\sigma}\le x^{-1/2}$ on $Q$. So the norm of this part
is at most $(0.02\pi a_k^{1/2}/(\Lambda^2\Gamma))^{1/2}\le(X_0V^3a_k^{\theta-1/2})^{1/2}\le(X_0V^3)^{1/2}a_k^{-5\cdot10^{-5}}$.
\item \emph{The Dirichlet sum} $S:=\sum_{n\in\mathcal B,\ \Xm\le n<a_k}n^{-s}$. For $\Xm\le u\le v\le a_k$, the elements of $\mathcal B$ in
$[u,v]$ number $D(v)-D(u-)+\int_{[u,v]\cap[X_0,\infty)}T'_{\mathrm{cont}}$, where $T'_{\mathrm{cont}}\le1+|R'|\le1.01$ and $|D|\le x^{\eta}$.
So a window of length $2$ in $[\Xm,2N+2]$ holds at most $K(N):=2(2N)^{\eta}+5$ of them, and $(N,2N]$ holds at most $1.02N$. Split $S$
into blocks $S_N$ over $n\in(N,2N]$, $N=2^{j-1}\Xm$ ($j\ge0$). List $\mathcal B\cap(N,2N]$ in increasing order, and split it into
$\lfloor K(N)\rfloor$ classes by the index modulo $\lfloor K(N)\rfloor$. Consecutive elements $n<n'$ of a class are $\lfloor K(N)\rfloor$ places
apart, so $n'>n+2$ and $\log n'-\log n\ge\log(1+1/N)\ge1/(N+1)$. By \eqref{eq:MV} for each class, and Cauchy--Schwarz over the classes,
\[
  \int_{\Gamma-1/4}^{2\Gamma+1/4}|S_N(\sigma+it)|^2\,dt\le K(N)\bigl(\Gamma+\tfrac12+3\pi(N+1)\bigr)\cdot1.02N^{1-2\sigma}.
\]
For $\sigma\ge\frac34$, $N\ge\Xm/2$ and $\eta\le\frac18$ this gives $\|S_N\|^2\le1.2N^{\eta-1/2}(1+10N/\Gamma)$. Summing $\|S_N\|$ over the
blocks gives two geometric series, with ratios $\le2^{-3/16}$ and $\ge2^{1/4}$; so, by \eqref{eq:Gammalow},
\[
  \|S\|\le11X_0^{-3/16}+22\,a_k^{\eta/2+1/4}\,\Gamma^{-1/2}\le11X_0^{-3/16}+32(X_0V^3)^{1/2}a_k^{(\theta+\eta)/2-1/4}.
\]
Here $(\theta+\eta)/2-\frac14\le-2.5\cdot10^{-5}$, since $\eta\le(\frac12-\theta)/2$ and $\theta\le\frac12-10^{-4}$.
\item \emph{Conclusion.} Since $V\ge X_0^2$ and $X_0\ge e^{10^8}$, each of the other four bounds is below $0.1X_0^{-3/16}$. So
$\|f\|\le12X_0^{-3/16}$.
\end{itemize}
Intersecting with (i), which holds on a set of measure $\ge0.036\,\Gamma$, leaves $\mathcal S_k$ with measure $\ge0.018\,\Gamma$ on
$\Good(<a_k+H(a_k))$.
\end{proof}

% Logged correction 64: the mean value theorem with an unspecified constant is now only a remark; the proof does not use it.
\begin{remark}
The mean value theorem for generalised Dirichlet polynomials \cite[Thm.~2.1]{Broucke2024} (from \cite{BrouckeDebruyne2023}) gives the
bound $\int_\Gamma^{2\Gamma}|S_N|^2\ll(\Gamma+N)K(N)N^{1-2\sigma}$ directly, with an implied constant that is not explicit. The split into
well-spaced classes above makes the constant explicit, so that (ii) holds for every admissible $X_0$.
\end{remark}

\begin{theorem}[Zeros]\label{thm:zeros}
% v3: Theorem 2.4
% Logged correction 76: gamma_k is the height defined by the rule of Lemma goodheights (for every omega).
Let $\gamma_k$ be as in Lemma~\ref{lem:goodheights}, and $\alpha_k:=\kappa_0x_k^{1-\theta}/(\gamma_k\log^3x_k)\in[1,2]$. On
$\mathcal G$, $\zeta_{\PP}$ has a zero $\rho_k$ with $|\rho_k-s^*_k|\le1/\log x_k$ for every $k$. Hence
\[
  \operatorname{Re}\rho_k=1-\frac{3\log\log x_k+O_{X_0}(1)}{\log x_k}\longrightarrow1.
\]
\end{theorem}
\begin{proof}
Apply Rouch\'e's theorem on the circle of Lemma~\ref{lem:designedzero}, comparing $\zeta_{\PP}$ with $Z_B+M_k$. The perturbation is
bounded by
% Logged correction 67: the two o(1) terms of v3 are replaced by explicit bounds.
\[
  |E_k|+\sum_{j\ne k}|\widehat R_j|+|\widehat D_{<a_k}|+|\widehat D_{\ge a_k}|+|\widehat E_{\Abs}|
  \le3e^{-0.49X_0^2}+\tfrac1{16}+165\Lambda e^{-X_0^2/8}+e^{-10^{10}}<\tfrac14,
\]
by the bullets after \eqref{eq:Rhat} (the circle lies in $0.99\le\sigma\le2$, $|t-\gamma_k|\le1$), Lemmas~\ref{lem:goodheights},
\ref{lem:latenoise} and~\ref{lem:ZA}, and $x_k\ge x_1\ge e^{X_0^2}$; here $\gamma_k\in\mathcal S_k$ on $\mathcal G$, and the circle lies in the
disc $|s-(1+i\gamma_k)|\le\delta_k\le\frac18$ (step~6 of the proof of Lemma~\ref{lem:designedzero}). The hypotheses of
Lemma~\ref{lem:designedzero} hold: on the disc, $|Z_B-Z_B(1+i\gamma_k)|\le\delta_k\cdot4.01\Lambda^2\le0.01$ (step~3 of that
proof), so $|Z_B|\ge0.59$ by Lemma~\ref{lem:goodheights}(i); and $|Z_B|\le Z_{\max}$ by Lemma~\ref{lem:ZB}.
\end{proof}

\begin{remark}
$1-\operatorname{Re}\rho_k\asymp3(1-\theta)\log\log\gamma_k/\log\gamma_k$. This is outside Landau's region
$1-c(1-\theta)/\log|t|$ by a factor $\log\log\gamma_k$, so it does not contradict it. In this design the loss comes from the
amplitude condition $\kappa'_k=\kappa_0/\log^2x_k$, which gives $3\log\log$ rather than $1\log\log$. Positivity \eqref{eq:G1}
alone would allow a constant $\kappa'_k=\kappa_0$. With the estimate used in Step~4 of Proposition~\ref{prop:g1}, which
bounds the oscillation response by the weight $\sum\mu^2n^{-\theta}$, the margin needs a decaying amplitude: the response of the drift to the
oscillation $R_k$ must stay $o(1/\log x)$ on $x_k^2<x<\exp(C'\log^2x_k)$, $C'=(4(1-\theta)/c_{\mathrm{LP}})^2$, where $\gamma_k$ is
not yet resolvable (proof of Proposition~\ref{prop:g1}, Step~4), and this requires $\kappa'_k\log x_k\to0$. Any power
$\log^{-1-\delta}x_k$ with $\delta>0$ would do, giving $(2+\delta)\log\log$; the exponent $2$ is a choice of the design.
% Logged correction 111 (CORRECTIONS-vs-v3.md): the need for a decaying amplitude is a property of the Step 4 estimate, not of the design.
A sharper estimate of the response (for instance through the Saias weights of Corollary~\ref{cor:saias}) may allow a
constant amplitude and hence $1\log\log$; we have not carried this out.
\end{remark}
\par\endgroup
